\documentclass[aspectratio=169]{beamer}

% Tema UFS --- Universidade Federal de Sergipe
\usetheme{UFS}

% Pacotes base
\usepackage[utf8]{inputenc}
\usepackage[portuguese]{babel}
\usepackage{amsmath,amssymb}
\usepackage{graphicx}
\usepackage{booktabs}
\usepackage{xcolor}

% TikZ e bibliotecas
\usepackage{tikz}
\usetikzlibrary{shapes.geometric, shapes.multipart, arrows.meta, positioning,
  matrix, fit, backgrounds, decorations.pathreplacing, patterns, calc}

% Listagens --- paleta VS Code Light+
\usepackage{listings}
\definecolor{bgcolor}{HTML}{FFFFFF}
\definecolor{linecolor}{HTML}{DDDDDD}
\definecolor{numcolor}{HTML}{999999}
\definecolor{kwcolor}{HTML}{0000FF}
\definecolor{strcolor}{HTML}{A31515}
\definecolor{cmtcolor}{HTML}{008000}

\lstdefinestyle{ufsStyle}{
  backgroundcolor=\color{bgcolor},
  basicstyle=\ttfamily\footnotesize\color{black},
  keywordstyle=\color{kwcolor}\bfseries,
  stringstyle=\color{strcolor},
  commentstyle=\color{cmtcolor}\itshape,
  numberstyle=\tiny\color{numcolor},
  numbers=left, numbersep=12pt,
  xleftmargin=20pt, framexleftmargin=20pt,
  breaklines=true, tabsize=4,
  keepspaces=true, showspaces=false,
  showstringspaces=false, showtabs=false,
  frame=single, rulecolor=\color{linecolor},
  framesep=4pt, captionpos=b, upquote=true,
  columns=fullflexible,
}
\lstset{style=ufsStyle, language=C}

\lstdefinestyle{tinyC}{
  style=ufsStyle, language=C, numbers=none,
  basicstyle=\ttfamily\scriptsize\color{black},
}
\lstdefinestyle{codC}{
  style=tinyC, basicstyle=\ttfamily\fontsize{7.5}{8.6}\selectfont\color{black},
}
\lstdefinestyle{shellStyle}{
  style=ufsStyle, language=bash, numbers=none,
  basicstyle=\ttfamily\scriptsize\color{black},
}

% --- Estilos TikZ reutilizados ---
\tikzset{
  cel/.style={rectangle, draw=UFSBlue, fill=UFSBlue!8, minimum width=0.55cm,
    minimum height=0.42cm, inner sep=0pt, font=\scriptsize\ttfamily},
  celvazia/.style={cel, fill=white, draw=UFSGray!60, text=UFSGray},
  celmorta/.style={cel, fill=black!5, draw=UFSGray!60, text=UFSGray!70},
  celr/.style={cel, draw=UFSRed, fill=UFSRed!10},
  var/.style={rectangle, draw=UFSGray, fill=black!5, minimum width=0.7cm,
    minimum height=0.5cm, font=\scriptsize\ttfamily},
  leg/.style={font=\scriptsize, text=UFSGray},
  seta/.style={-{Stealth[length=1.8mm]}, thick, UFSBlue},
  setar/.style={-{Stealth[length=1.8mm]}, thick, UFSRed},
  setag/.style={-{Stealth[length=1.8mm]}, thick, UFSGray},
}

% Fila desenhada da esquerda (frente) para a direita (fim):
% \fila{x}{y}{{a},{b},{c}}
\newcommand{\fila}[3]{%
  \foreach \e [count=\k from 0] in {#3} {
    \node[cel] at (#1+0.55*\k, #2) {\e};
  }
}

% Referencia da fonte no rodape do frame
\newcommand{\fonte}[1]{\vfill\hfill{\scriptsize\textcolor{UFSGray}{Fonte: #1}}}

% --- Metadados ---
\title{Filas: conceito e implementação}
\subtitle{Entra-se por uma ponta e sai-se pela outra}
\author{Prof.\ André Yoshiaki Kashiwabara}
\institute{DCOMP -- Universidade Federal de Sergipe\\
\texttt{andreyoshiaki@dcomp.ufs.br}}
\date{}

\begin{document}

\begin{frame}[plain]
  \titlepage
\end{frame}

\begin{frame}{O que você vai aprender hoje}
  \begin{center}
  \begin{tikzpicture}[
    item/.style={rectangle, draw=UFSBlue, fill=UFSBlue!10, rounded corners,
      minimum width=10.4cm, minimum height=0.5cm, font=\small}
  ]
    \node[item] (t1) at (0, 0)    {Fila: o primeiro a entrar é o primeiro a sair};
    \node[item] (t2) at (0,-0.64) {Fila num vetor: \texttt{ini} e \texttt{fim} só andam para a direita};
    \node[item] (t3) at (0,-1.28) {A fila circular: o fim dá a volta no vetor};
    \node[item] (t4) at (0,-1.92) {Quando o vetor circular enche: dobrar sem embaralhar};
    \node[item] (t5) at (0,-2.56) {Fila numa lista: sai do primeiro nó, entra depois do último};
    \node[item] (t6) at (0,-3.20) {Quanto custa: deslocar, dar a volta ou encadear};
    \node[item] (t7) at (0,-3.84) {Exercícios};
    \draw[seta] (t1) -- (t2);
    \draw[seta] (t2) -- (t3);
    \draw[seta] (t3) -- (t4);
    \draw[seta] (t4) -- (t5);
    \draw[seta] (t5) -- (t6);
    \draw[seta] (t6) -- (t7);
  \end{tikzpicture}
  \end{center}
\end{frame}

% ==================================================================
\section{A pergunta de hoje}
% ==================================================================

\begin{frame}{O problema: quem fica com a batata?}
  \begin{columns}[T]
    \begin{column}{0.5\textwidth}
      \textbf{Até aqui}\\[0.3em]
      \small
      A pilha devolve tudo na ordem \emph{inversa}: o último a entrar sai
      primeiro. Serviu para parênteses, desfazer e a pilha de execução.
    \end{column}
    \begin{column}{0.5\textwidth}
      \textbf{O que vamos descobrir}\\[0.3em]
      \small
      Muitos problemas querem a ordem de \textbf{chegada}: quem chegou antes
      é atendido antes. A estrutura é tão simples quanto a pilha --- mas
      guardá-la num vetor esconde uma armadilha.
    \end{column}
  \end{columns}
  \begin{block}{A pergunta de hoje}
    $n$ crianças em roda, numeradas de 1 a $n$. A batata quente passa de mão
    em mão $k$ vezes; quem fica com ela sai, e a brincadeira continua a
    partir do seguinte. Quem sobra?
  \end{block}
\end{frame}

\begin{frame}[fragile]{O programa de hoje: \texttt{batata.c} (1/2)}
  \begin{columns}[T]
    \begin{column}{0.46\textwidth}
\begin{lstlisting}[style=codC]
#include <stdio.h>
#include <stdlib.h>
#define MAX 100
typedef struct {
    int v[MAX]; /* frente: v[ini] */
    int ini, fim; /* entra em v[fim] */
} Fila;

void enfileira(Fila *f, int x)
{
    if (f->fim == MAX) {
        fprintf(stderr, "fila cheia\n");
        exit(1);
    }
    f->v[f->fim] = x;
    f->fim++;
}
\end{lstlisting}
    \end{column}
    \begin{column}{0.5\textwidth}
\begin{lstlisting}[style=codC]
int desenfileira(Fila *f)
{
    if (f->ini == f->fim) {
        fprintf(stderr, "fila vazia\n");
        exit(1);
    }
    int x = f->v[f->ini];
    f->ini++;
    return x;
}

int tamanho(const Fila *f)
{
    return f->fim - f->ini;
}
\end{lstlisting}
    \end{column}
  \end{columns}
\end{frame}

\begin{frame}[fragile]{\texttt{batata.c} (2/2): a roda e o \texttt{main}}
  \vspace{-0.4em}
  \begin{columns}[T]
    \begin{column}{0.52\textwidth}
\begin{lstlisting}[style=codC]
/* imprime quem sai; devolve quem fica */
int roda(int n, int k)
{
    Fila f = {.ini = 0, .fim = 0};
    for (int i = 1; i <= n; i++)
        enfileira(&f, i);
    while (tamanho(&f) > 1) {
        for (int j = 0; j < k; j++)
            /* passa a batata */
            enfileira(&f, desenfileira(&f));
        /* quem ficou com ela sai */
        printf(" %d", desenfileira(&f));
    }
    return desenfileira(&f);
}
\end{lstlisting}
    \end{column}
    \begin{column}{0.46\textwidth}
\begin{lstlisting}[style=codC]
int main(void)
{
    int casos[][2] = {
        {6, 0}, {5, 2}, {7, 1},
    };
    for (int i = 0; i < 3; i++) {
        printf("n = %d, k = %d: sai",
               casos[i][0], casos[i][1]);
        int fica = roda(casos[i][0],
                        casos[i][1]);
        printf(" | fica %d\n", fica);
    }
    return 0;
}
\end{lstlisting}
    \end{column}
  \end{columns}
\end{frame}

\begin{frame}{O que este programa imprime?}
  \begin{block}{Antes de compilar, escreva sua resposta}
    \begin{enumerate}
      \item Para $k = 0$, a batata nem passa. Em que ordem as seis crianças
        saem?
      \item Para $n = 5$ e $k = 2$, quem sai primeiro? Quem sobra?
      \item E para $n = 7$ e $k = 1$?
      \item Passar a batata é \texttt{enfileira(\&f, desenfileira(\&f))}. O
        que essa linha faz com a criança da frente?
    \end{enumerate}
  \end{block}
\end{frame}

\begin{frame}[fragile]{Compile, execute e compare}
\begin{lstlisting}[style=shellStyle]
$ gcc -Wall -Wextra batata.c -o batata
$ ./batata
n = 6, k = 0: sai 1 2 3 4 5 | fica 6
n = 5, k = 2: sai 3 1 5 2 | fica 4
n = 7, k = 1: sai 2 4 6 1 5 3 | fica 7
\end{lstlisting}
  \begin{columns}[T]
    \begin{column}{0.48\textwidth}
      \begin{block}{O que o programa guarda}
        \small A roda inteira, na ordem: quem está com a batata na
        \textbf{frente}, quem acabou de passá-la no \textbf{fim}.
      \end{block}
    \end{column}
    \begin{column}{0.48\textwidth}
      \small Passar a batata tira a criança da frente e a põe no fim: a roda
      gira sem que ninguém mude de lugar em relação aos outros. Com
      $k = 0$, saem na ordem em que entraram.
    \end{column}
  \end{columns}
  \fonte{Graham, Knuth e Patashnik (1994), §1.3, \emph{o problema de Josephus}~\cite{graham1994concrete}}
\end{frame}

% ==================================================================
\section{O que é uma fila}
% ==================================================================

\begin{frame}{Definição: fila}
  \begin{columns}[c]
    \begin{column}{0.52\textwidth}
      \begin{block}{Fila}
        \small Uma sequência em que as inserções acontecem numa ponta, o
        \emph{fim}, e as remoções na outra, a \emph{frente}. O primeiro a
        entrar é o primeiro a sair: \emph{FIFO}, \emph{first in, first out}.
      \end{block}
      \small
      \begin{tabular}{@{}ll@{}}
        \texttt{enfileira(f, x)} & põe \texttt{x} no fim\\
        \texttt{desenfileira(f)} & tira e devolve a frente\\
        \texttt{frente(f)}       & devolve a frente sem tirar\\
        \texttt{vazia(f)}        & a fila não tem nada?
      \end{tabular}
    \end{column}
    \begin{column}{0.44\textwidth}
      \begin{center}
      \begin{tikzpicture}
        \fila{0}{0}{{3},{4},{5},{1},{2}}
        \node[leg] at (0,-0.45) {frente};
        \node[leg] at (2.2,-0.45) {fim};
        \draw[seta] (-0.3,0) -- (-1.0,0) node[left, leg, text=UFSBlue] {sai};
        \node[celr] (novo) at (3.4,0) {6};
        \draw[setar] (novo.west) -- (2.5,0);
        \node[leg, text=UFSRed] at (3.4,0.45) {entra};
      \end{tikzpicture}
      \end{center}
      \small Na pilha, as duas operações usam a mesma ponta. Na fila, cada
      uma usa a sua --- é daí que vem toda a diferença na implementação.
    \end{column}
  \end{columns}
  \fonte{Cormen et al.\ (2009), §10.1~\cite{cormen2009introduction};
    Celes et al.\ (2004), cap.~12~\cite{celes2004introducao}}
\end{frame}

\begin{frame}{Seguindo a fila: $n = 5$, $k = 2$}
  \begin{center}
  \begin{tikzpicture}
    \node[leg, anchor=east] at (-0.5,0) {início};
    \fila{0}{0}{{1},{2},{3},{4},{5}}

    \node[leg, anchor=east] at (-0.5,-0.6) {passa 2};
    \fila{0}{-0.6}{{3},{4},{5},{1},{2}}
    \node[leg, anchor=west, text=UFSRed] at (2.7,-0.6) {sai \texttt{3}};

    \node[leg, anchor=east] at (-0.5,-1.2) {passa 2};
    \fila{0}{-1.2}{{1},{2},{4},{5}}
    \node[leg, anchor=west, text=UFSRed] at (2.7,-1.2) {sai \texttt{1}};

    \node[leg, anchor=east] at (-0.5,-1.8) {passa 2};
    \fila{0}{-1.8}{{5},{2},{4}}
    \node[leg, anchor=west, text=UFSRed] at (2.7,-1.8) {sai \texttt{5}};

    \node[leg, anchor=east] at (-0.5,-2.4) {passa 2};
    \fila{0}{-2.4}{{2},{4}}
    \node[leg, anchor=west, text=UFSRed] at (2.7,-2.4) {sai \texttt{2}};
    \node[leg, anchor=west] at (2.7,-3.0) {fica \texttt{4}};
  \end{tikzpicture}
  \end{center}
  \small Cada passada tira da frente e põe no fim; a criança que sai é a que
  está na frente depois das $k$ passadas. Ninguém precisa ser procurado.
  \fonte{Graham, Knuth e Patashnik (1994), §1.3~\cite{graham1994concrete};
    Tenenbaum et al.\ (1995), cap.~4~\cite{tenenbaum1995estruturas}}
\end{frame}

\begin{frame}{Filas que você já usou}
  \begin{columns}[T]
    \begin{column}{0.32\textwidth}
      \begin{block}{O teclado}
        \small O que você digita espera num \emph{buffer} até o
        \texttt{scanf} ler --- na ordem em que foi digitado. É por isso que
        o \texttt{\textbackslash n} esquecido aparece no próximo
        \texttt{scanf}.
      \end{block}
    \end{column}
    \begin{column}{0.32\textwidth}
      \begin{block}{A impressora}
        \small Os documentos esperam na fila de impressão; o primeiro
        mandado é o primeiro impresso.
      \end{block}
    \end{column}
    \begin{column}{0.32\textwidth}
      \begin{block}{O processador}
        \small O sistema dá a cada programa uma fatia de tempo e o põe no fim
        da fila de prontos --- é a batata quente, sem ninguém sair.
      \end{block}
    \end{column}
  \end{columns}
  \medskip
  \small Em todos, a pergunta é a mesma: \emph{quem está esperando há mais
  tempo?}
  \fonte{Tanenbaum e Bos (2015), §2.4.3, \emph{round-robin}~\cite{tanenbaum2015modern};
    Kernighan e Ritchie (1988), §8.5~\cite{kernighan1988c}}
\end{frame}

\begin{frame}{Recapitulando: o que é uma fila}
  \begin{block}{O que aprendemos}
    \begin{itemize}
      \item fila: entra no fim, sai da frente --- o primeiro a entrar é o
        primeiro a sair;
      \item quatro operações: \texttt{enfileira}, \texttt{desenfileira},
        \texttt{frente} e \texttt{vazia};
      \item passar a batata é desenfileirar e enfileirar de novo: a fila
        ``gira'' sem mudar a ordem relativa;
      \item buffers, impressoras e o escalonador do sistema são filas.
    \end{itemize}
  \end{block}
\end{frame}

% ==================================================================
\section{Fila num vetor}
% ==================================================================

\begin{frame}{A frente é \texttt{v[ini]}, o próximo entra em \texttt{v[fim]}}
  \begin{center}
  \begin{tikzpicture}
    \foreach \k in {0,1} {
      \node[celmorta, minimum width=0.7cm] (v\k) at (0.7*\k,0) {};
      \node[leg] at (0.7*\k,-0.4) {\k};
    }
    \foreach \c [count=\k from 2] in {3,4,5,1,2} {
      \node[cel, minimum width=0.7cm] (v\k) at (0.7*\k,0) {\c};
      \node[leg] at (0.7*\k,-0.4) {\k};
    }
    \foreach \k in {7,8} {
      \node[celvazia, minimum width=0.7cm] (v\k) at (0.7*\k,0) {?};
      \node[leg] at (0.7*\k,-0.4) {\k};
    }
    \node[leg, anchor=west] at (6.0,0) {\ldots\ \texttt{MAX-1}};
    \draw[setar] (v2.north) ++(0,0.5) node[above, leg, text=UFSRed] {frente: \texttt{v[ini]}, \texttt{ini = 2}} -- (v2.north);
    \draw[seta] (v7.north) ++(0,0.5) node[above, leg, text=UFSBlue] {entra: \texttt{v[fim]}, \texttt{fim = 7}} -- (v7.north);
  \end{tikzpicture}
  \end{center}
  \begin{columns}[T]
    \begin{column}{0.48\textwidth}
      \small
      \texttt{enfileira}: escreve em \texttt{v[fim]}, depois \texttt{fim++}.\\
      \texttt{desenfileira}: lê \texttt{v[ini]}, depois \texttt{ini++}.\\
      \texttt{vazia}: \texttt{ini == fim}; \texttt{tamanho}: \texttt{fim - ini}.
    \end{column}
    \begin{column}{0.48\textwidth}
      \small Nenhuma operação tem laço. Mas repare: na pilha só o topo
      andava, e para os dois lados. Aqui \texttt{ini} e \texttt{fim}
      \textbf{só aumentam}. As posições à esquerda de \texttt{ini} nunca
      mais são usadas.
    \end{column}
  \end{columns}
  \fonte{Cormen et al.\ (2009), §10.1, \textsc{Enqueue} e \textsc{Dequeue}~\cite{cormen2009introduction};
    Celes et al.\ (2004), cap.~12~\cite{celes2004introducao}}
\end{frame}

\begin{frame}[fragile]{E se a roda for maior?}
  \begin{columns}[c]
    \begin{column}{0.5\textwidth}
\begin{lstlisting}[style=shellStyle]
$ ./batata 41 2
\end{lstlisting}
      \small O \texttt{main} completo aceita \texttt{n} e \texttt{k} pela
      linha de comando. São 41 crianças --- a fila nunca tem mais do que 41
      elementos --- e \texttt{MAX} é 100.
    \end{column}
    \begin{column}{0.46\textwidth}
      \begin{block}{Pense antes de passar o slide}
        \small
        \begin{enumerate}
          \item Cabe? Algum dos dois testes de \texttt{enfileira} e
            \texttt{desenfileira} vai disparar?
          \item Quantas vezes \texttt{enfileira} é chamada até sobrar uma
            criança?
          \item Quanto vale \texttt{fim} no final?
        \end{enumerate}
      \end{block}
    \end{column}
  \end{columns}
\end{frame}

\begin{frame}[fragile]{Execute e compare}
\begin{lstlisting}[style=shellStyle]
$ ./batata 41 2
n = 41, k = 2: sai 3 6 9 12 15 18 21 24 27 30 33 36 39 1 5 10 14 19 23 28
                   32 37 41 7 13 20 26 34 40
enfileira: fila cheia
\end{lstlisting}
  \begin{columns}[T]
    \begin{column}{0.48\textwidth}
      \small Saíram 29 crianças; restavam 12 na roda. A fila ``encheu''
      com 12 elementos num vetor de 100.
    \end{column}
    \begin{column}{0.48\textwidth}
      \small Desta vez não houve lixo nem estouro silencioso: o teste de
      \texttt{enfileira} --- a lição da aula de pilhas --- parou o programa.
      Mas parou um programa que deveria funcionar.
    \end{column}
  \end{columns}
\end{frame}

\begin{frame}{A resposta: a fila anda para a direita}
  \begin{center}
  \begin{tikzpicture}
    \node[celmorta, minimum width=0.7cm] at (0,0) {};
    \node[leg] at (0,-0.4) {0};
    \node[leg] at (0.7,0) {\ldots};
    \node[celmorta, minimum width=0.7cm] at (1.4,0) {};
    \node[leg] at (1.4,-0.4) {88};
    \foreach \c [count=\k from 0] in {8,11,16} {
      \node[cel, minimum width=0.7cm] (o\k) at (2.1+0.7*\k,0) {\c};
    }
    \node[leg] at (2.1,-0.4) {89};
    \node[leg] at (4.55,0) {\ldots};
    \node[cel, minimum width=0.7cm] at (5.3,0) {2};
    \node[leg] at (5.3,-0.4) {99};
    \node[celr, minimum width=0.7cm] (fora) at (6.2,0) {};
    \node[leg, text=UFSRed] at (6.2,-0.4) {100};
    \draw[UFSGray, thick, decorate, decoration={brace, amplitude=4pt}]
      (-0.35,0.3) -- (1.75,0.3) node[midway, above=3pt, leg] {89 posições abandonadas};
    \draw[UFSBlue, thick, decorate, decoration={brace, amplitude=4pt}]
      (1.75,0.3) -- (5.65,0.3) node[midway, above=3pt, leg, text=UFSBlue] {11 na fila};
    \draw[setar] (6.2,1.1) node[above, leg, text=UFSRed] {\texttt{fim == MAX}: o 4 não entra} -- (fora.north);
  \end{tikzpicture}
  \end{center}
  \small
  Cada criança que entra e cada passada da batata é um \texttt{enfileira}:
  são $n + k(n-1) = 41 + 2 \cdot 40 = 121$ ao todo, e \texttt{fim} só cresce. A fila nunca passou de
  41 elementos, mas foi andando para a direita e bateu na parede, deixando
  para trás posições vazias que ninguém mais usa.

  \medskip
  Um vetor de 100 só aguenta 100 \texttt{enfileira} \textbf{na vida inteira}
  da fila. Não basta saber quantos elementos cabem: é preciso reaproveitar o
  espaço que a frente deixou.
  \fonte{Cormen et al.\ (2009), §10.1~\cite{cormen2009introduction};
    Tenenbaum et al.\ (1995), cap.~4~\cite{tenenbaum1995estruturas}}
\end{frame}

\begin{frame}[fragile]{Primeira ideia: trazer todo mundo para a frente}
  \begin{columns}[c]
    \begin{column}{0.5\textwidth}
\begin{lstlisting}[style=codC]
int desenfileira(Fila *f)
{   /* a frente e sempre v[0] */
    int x = f->v[0];
    for (int i = 1; i < f->fim; i++)
        f->v[i - 1] = f->v[i];
    f->fim--;
    return x;
}
\end{lstlisting}
    \end{column}
    \begin{column}{0.46\textwidth}
      \small Sem \texttt{ini}: quem sai é sempre \texttt{v[0]}, e os outros
      andam uma casa para a esquerda. Agora \texttt{MAX} = 41 basta para a
      roda de 41.

      \medskip
      É assim numa fila de banco: quando o primeiro é atendido, todos dão
      um passo. Mas no vetor cada passo é uma cópia --- com $n$ na fila,
      \textbf{cada} \texttt{desenfileira} custa $n - 1$ cópias.
    \end{column}
  \end{columns}
  \fonte{Celes et al.\ (2004), cap.~12~\cite{celes2004introducao};
    medido com \texttt{custo.c}, adiante}
\end{frame}

\begin{frame}{Recapitulando: fila num vetor}
  \begin{block}{O que aprendemos}
    \begin{itemize}
      \item com \texttt{ini} e \texttt{fim}, enfileirar e desenfileirar não
        têm laço: escrevem em \texttt{v[fim]} e leem de \texttt{v[ini]};
      \item mas os dois índices só crescem: a fila anda pelo vetor e
        ``enche'' depois de \texttt{MAX} enfileiramentos, não com
        \texttt{MAX} elementos;
      \item deslocar todos a cada saída resolve o espaço, mas faz cada
        \texttt{desenfileira} custar o tamanho da fila;
      \item falta reaproveitar o espaço que a frente deixa para trás.
    \end{itemize}
  \end{block}
\end{frame}

% ==================================================================
\section{A fila circular}
% ==================================================================

\begin{frame}{O fim dá a volta}
  \begin{columns}[c]
    \begin{column}{0.42\textwidth}
      \begin{center}
      \begin{tikzpicture}
        \foreach \i/\c/\s in {0/7/cel, 1/8/cel, 2/{}/celvazia, 3/{}/celvazia,
                              4/{}/celvazia, 5/3/cel, 6/4/cel, 7/5/cel} {
          \node[\s, minimum width=0.55cm] (c\i) at ({90-45*\i}:1.3) {\c};
          \node[leg] at ({90-45*\i}:0.8) {\i};
        }
        \draw[setar] (c5) ++(-0.8,-0.5) node[left, leg, text=UFSRed] {ini} -- (c5);
        \draw[seta] (c2) ++(0.9,0) node[right, leg, text=UFSBlue] {fim} -- (c2);
        \draw[setag, -{Stealth[length=1.6mm]}] (100:1.85) arc (100:30:1.85);
      \end{tikzpicture}
      \end{center}
    \end{column}
    \begin{column}{0.54\textwidth}
      \small Imagine o vetor dobrado em anel: depois de \texttt{v[MAX-1]} vem
      \texttt{v[0]}. Quando o fim chega à última posição, continua no começo
      --- no espaço que a frente já liberou.

      \medskip
      Em C, ``dar a volta'' é o resto da divisão:
      \begin{center}
        \texttt{fim = (fim + 1) \% MAX;}\\
        \texttt{ini = (ini + 1) \% MAX;}
      \end{center}
      Na figura, \texttt{MAX = 8}: a fila é \texttt{3 4 5 7 8}, com a frente
      em \texttt{v[5]} e o próximo entrando em \texttt{v[2]}.
    \end{column}
  \end{columns}
  \fonte{Cormen et al.\ (2009), §10.1, Figura 10.2~\cite{cormen2009introduction};
    Knuth (1997), §2.2.2~\cite{knuth1997art}}
\end{frame}

\begin{frame}{Cheia ou vazia?}
  \begin{columns}[c]
    \begin{column}{0.5\textwidth}
      \begin{center}
      \begin{tikzpicture}
        \node[leg, anchor=east] at (-0.5,0) {vazia};
        \foreach \k in {0,...,3} \node[celvazia] at (0.55*\k,0) {};
        \node[leg, anchor=west] at (2.0,0) {\texttt{ini = 0, fim = 0}};
        \node[leg, anchor=east] at (-0.5,-0.7) {+ 4 entram};
        \fila{0}{-0.7}{{1},{2},{3},{4}}
        \node[leg, anchor=west] at (2.0,-0.7) {\texttt{ini = 0, fim = 0}};
      \end{tikzpicture}
      \end{center}
      \small \texttt{MAX = 4}. Depois do quarto \texttt{enfileira}, o
      \texttt{fim} deu a volta: \texttt{(3 + 1) \% 4 == 0}.
    \end{column}
    \begin{column}{0.46\textwidth}
      \begin{block}{Pense antes de passar o slide}
        \small
        \begin{enumerate}
          \item Com \texttt{ini} e \texttt{fim} apenas, como \texttt{vazia}
            distingue as duas situações?
          \item \texttt{tamanho} pode continuar sendo \texttt{fim - ini}?
          \item O que mais a fila precisaria guardar?
        \end{enumerate}
      \end{block}
    \end{column}
  \end{columns}
\end{frame}

\begin{frame}{A resposta: guarde quantos há}
  \begin{columns}[T]
    \begin{column}{0.48\textwidth}
      \small Com \texttt{ini} e \texttt{fim}, a fila cheia e a vazia são o
      mesmo estado: \texttt{ini == fim}. E \texttt{fim - ini} pode ser
      negativo quando o fim já deu a volta.

      \medskip
      Há duas saídas clássicas:
      \begin{itemize}
        \item guardar \texttt{n}, o número de elementos: vazia é
          \texttt{n == 0}, cheia é \texttt{n == MAX}, e o \texttt{fim}
          nem precisa ser guardado;
        \item deixar sempre uma posição vazia: cheia é
          \texttt{(fim + 1) \% MAX == ini}.
      \end{itemize}
    \end{column}
    \begin{column}{0.48\textwidth}
      \begin{block}{Com \texttt{ini} e \texttt{n}}
        \small
        \begin{tabular}{@{}ll@{}}
          frente      & \texttt{v[ini]}\\
          próximo     & \texttt{v[(ini + n) \% MAX]}\\
          vazia       & \texttt{n == 0}\\
          cheia       & \texttt{n == MAX}\\
        \end{tabular}
      \end{block}
      \small O livro de Cormen usa a segunda saída e guarda no máximo
      $\texttt{MAX} - 1$ elementos. Aqui usaremos a primeira.
    \end{column}
  \end{columns}
  \fonte{Cormen et al.\ (2009), §10.1~\cite{cormen2009introduction};
    Knuth (1997), §2.2.2~\cite{knuth1997art}}
\end{frame}

\begin{frame}[fragile]{A roda com a fila circular: \texttt{batata\_circular.c}}
  \begin{columns}[T]
    \begin{column}{0.5\textwidth}
\begin{lstlisting}[style=codC]
#define MAX 41

typedef struct {
    int v[MAX];
    int ini, n; /* frente, quantos */
} Fila;

void enfileira(Fila *f, int x)
{
    if (f->n == MAX) { /* cheia */ }
    f->v[(f->ini + f->n) % MAX] = x;
    f->n++;
}
\end{lstlisting}
    \end{column}
    \begin{column}{0.46\textwidth}
\begin{lstlisting}[style=codC]
int desenfileira(Fila *f)
{
    if (f->n == 0) { /* vazia */ }
    int x = f->v[f->ini];
    f->ini = (f->ini + 1) % MAX;
    f->n--;
    return x;
}
\end{lstlisting}
      \small \texttt{roda} e \texttt{main} não mudaram uma linha: só a
      fila por dentro. Os testes de cheia e vazia são os de
      \texttt{batata.c}.
    \end{column}
  \end{columns}
\end{frame}

\begin{frame}[fragile]{Execute e compare}
\begin{lstlisting}[style=shellStyle]
$ gcc -Wall -Wextra batata_circular.c -o batata_circular
$ ./batata_circular 41 2
n = 41, k = 2: sai 3 6 9 12 15 18 21 24 27 30 33 36 39 1 5 10 14 19 23 28
                   32 37 41 7 13 20 26 34 40 8 17 29 38 11 25 2 22 4 35 16 | fica 31
\end{lstlisting}
  \begin{columns}[T]
    \begin{column}{0.48\textwidth}
      \small Os 121 \texttt{enfileira} cabem num vetor de \textbf{41}: o
      espaço que a frente libera é reaproveitado pelo fim.
    \end{column}
    \begin{column}{0.48\textwidth}
      \small Conta a lenda que Josephus e um amigo, numa roda de 41 em que
      cada terceiro saía, escolheram os lugares 16 e 31 --- os dois
      últimos.
    \end{column}
  \end{columns}
  \fonte{Graham, Knuth e Patashnik (1994), §1.3~\cite{graham1994concrete}}
\end{frame}

\begin{frame}{Recapitulando: a fila circular}
  \begin{block}{O que aprendemos}
    \begin{itemize}
      \item o vetor vira anel com \texttt{\% MAX}: o fim continua no espaço
        que a frente liberou;
      \item com só \texttt{ini} e \texttt{fim}, cheia e vazia são o mesmo
        estado --- guarde \texttt{n} (ou deixe uma posição sempre livre);
      \item a próxima posição livre é \texttt{(ini + n) \% MAX};
      \item agora \texttt{MAX} limita quantos estão na fila, não quantos já
        passaram por ela.
    \end{itemize}
  \end{block}
\end{frame}

% ==================================================================
\section{Quando o vetor circular enche}
% ==================================================================

\begin{frame}[fragile]{Dobrar, como na pilha: \texttt{dobra\_errada.c}}
  \begin{columns}[T]
    \begin{column}{0.5\textwidth}
\begin{lstlisting}[style=codC]
void enfileira(Fila *f, int x)
{
    if (f->n == f->cap) { /* dobra */
        f->cap *= 2;
        f->v = realloc(f->v,
                 f->cap * sizeof *f->v);
    }
    f->v[(f->ini + f->n) % f->cap] = x;
    f->n++;
}
\end{lstlisting}
      \small Na pilha, crescer era um \texttt{realloc}. O \texttt{main}
      começa com \texttt{cap = 4}, enfileira de 1 a 4, desenfileira dois,
      enfileira 5, 6 e 7 --- o 7 faz dobrar --- e imprime tudo o que sair.
    \end{column}
    \begin{column}{0.46\textwidth}
      \begin{block}{Pense antes de passar o slide}
        \small
        \begin{enumerate}
          \item Desenhe o vetor de 4 antes do 7: onde estão \texttt{3 4 5 6}
            e quanto vale \texttt{ini}?
          \item Desenhe o vetor de 8 depois do \texttt{realloc}. Onde o 7
            é escrito?
          \item Em que ordem os cinco saem?
        \end{enumerate}
      \end{block}
    \end{column}
  \end{columns}
  \fonte{ISO/IEC 9899:2018, §7.22.3.5~\cite{iso2018c17}}
\end{frame}

\begin{frame}[fragile]{Execute e compare}
\begin{lstlisting}[style=shellStyle]
$ ./dobra_errada
cap = 8, ini = 2, n = 5
sai: 3 4 0 0 7
$ gcc -g -fsanitize=address dobra_errada.c -o dobra_errada-asan
$ ./dobra_errada-asan
cap = 8, ini = 2, n = 5
sai: 3 4 -1094795586 -1094795586 7
\end{lstlisting}
  \begin{columns}[T]
    \begin{column}{0.48\textwidth}
      \small O 5 e o 6 sumiram; no lugar deles saiu o que havia na memória
      nova --- zero numa execução, outro valor na outra.
    \end{column}
    \begin{column}{0.48\textwidth}
      \small O \emph{sanitizer} não acusa nada: todas as posições lidas
      estão \emph{dentro} do bloco. Ele só preenche a memória nova com um
      padrão, \texttt{0xBEBEBEBE}, e o padrão aparece na saída.
    \end{column}
  \end{columns}
\end{frame}

\begin{frame}{A resposta: o \texttt{realloc} copia o vetor, não a fila}
  \begin{center}
  \begin{tikzpicture}
    \node[leg, anchor=east] at (-0.5,0.6) {antes};
    \foreach \c [count=\k from 0] in {5,6,3,4}
      \node[cel] at (0.55*\k,0.6) {\c};
    \draw[setar] (1.1,1.3) node[above, leg, text=UFSRed] {\texttt{ini = 2}} -- (1.1,0.85);
    \node[leg, anchor=west] at (2.3,0.6) {fila: \texttt{3 4 5 6} (o 5 e o 6 deram a volta)};

    \node[leg, anchor=east] at (-0.5,-0.4) {depois};
    \foreach \c [count=\k from 0] in {5,6,3,4}
      \node[cel] at (0.55*\k,-0.4) {\c};
    \node[celvazia] at (2.2,-0.4) {?};
    \node[celvazia] at (2.75,-0.4) {?};
    \node[celr] at (3.3,-0.4) {7};
    \node[celvazia] at (3.85,-0.4) {?};
    \draw[setag] (1.1,-0.85) -- (1.1,-0.65);
    \node[leg, anchor=west] at (4.3,-0.4) {de \texttt{v[2]} em diante: \texttt{3 4 ? ? 7}};
  \end{tikzpicture}
  \end{center}
  \small
  O \texttt{realloc} copia as posições 0 a 3 para as posições 0 a 3 do bloco
  novo. Só que, na fila, o 5 e o 6 vinham \emph{depois} do 4: com
  \texttt{cap = 8}, a posição seguinte a \texttt{v[3]} deixou de ser
  \texttt{v[0]} e passou a ser \texttt{v[4]}, que é lixo.

  \medskip
  A conta \texttt{(ini + i) \% cap} mudou de significado quando \texttt{cap}
  mudou. Para crescer, é preciso copiar \textbf{na ordem da fila}, não na
  ordem do vetor.
  \fonte{ISO/IEC 9899:2018, §7.22.3.5~\cite{iso2018c17};
    Cormen et al.\ (2009), §17.4~\cite{cormen2009introduction}}
\end{frame}

\begin{frame}[fragile]{Crescer desenrolando o anel}
  \begin{columns}[T]
    \begin{column}{0.54\textwidth}
\begin{lstlisting}[style=codC]
/* fila_vetor.c: vetor com o dobro do
   tamanho, ja na ordem: frente em 0 */
static void cresce(Fila *f)
{
    int *novo = aloca(NULL,
                  2 * f->cap * sizeof *novo);
    for (int i = 0; i < f->n; i++)
        novo[i] = f->v[(f->ini + i) % f->cap];
    free(f->v);
    f->v = novo;
    f->ini = 0;
    f->cap *= 2;
}
\end{lstlisting}
    \end{column}
    \begin{column}{0.42\textwidth}
      \begin{center}
      \begin{tikzpicture}
        \foreach \c [count=\k from 0] in {3,4,5,6,7}
          \node[cel] at (0.55*\k,0) {\c};
        \foreach \k in {5,6,7}
          \node[celvazia] at (0.55*\k,0) {};
        \draw[setar] (0,0.75) node[above, leg, text=UFSRed] {\texttt{ini = 0}} -- (0,0.25);
      \end{tikzpicture}
      \end{center}
      \small O laço percorre a fila da frente para o fim, com o
      \texttt{cap} \emph{antigo}, e escreve a partir de \texttt{novo[0]}.

      \medskip
      Ainda são $n$ cópias por crescimento, como na pilha: dobrando, o custo
      por \texttt{enfileira} continua constante, amortizado.
    \end{column}
  \end{columns}
  \fonte{Cormen et al.\ (2009), §17.4, \emph{dynamic tables}~\cite{cormen2009introduction}}
\end{frame}

\begin{frame}{Recapitulando: quando o vetor circular enche}
  \begin{block}{O que aprendemos}
    \begin{itemize}
      \item numa fila que deu a volta, a ordem dos elementos não é a ordem
        das posições no vetor;
      \item o \texttt{realloc} preserva as posições, e a conta
        \texttt{\% cap} muda com o novo \texttt{cap}: a fila embaralha;
      \item cresça copiando na ordem da fila, com \texttt{ini} voltando a 0;
      \item dobrando, o custo amortizado de \texttt{enfileira} continua
        constante --- a mesma conta da pilha.
    \end{itemize}
  \end{block}
\end{frame}

% ==================================================================
\section{Fila numa lista}
% ==================================================================

\begin{frame}[fragile]{Sai do primeiro nó, entra depois do último: \texttt{fila\_lista.c}}
  \begin{columns}[T]
    \begin{column}{0.5\textwidth}
\begin{lstlisting}[style=codC]
struct fila { No *ini, *fim; };

void enfileira(Fila *f, int x)
{   /* insere no fim */
    No *novo = malloc(sizeof *novo);
    if (novo == NULL) exit(1);
    novo->valor = x; novo->prox = NULL;
    if (f->fim == NULL) f->ini = novo;
    else f->fim->prox = novo;
    f->fim = novo;
}
\end{lstlisting}
    \end{column}
    \begin{column}{0.46\textwidth}
\begin{lstlisting}[style=codC]
int desenfileira(Fila *f)
{   /* remove o primeiro */
    if (f->ini == NULL) {
        fprintf(stderr, "fila vazia\n");
        exit(1);
    }
    No *velho = f->ini;
    int x = velho->valor;
    f->ini = velho->prox;
    if (f->ini == NULL) f->fim = NULL;
    free(velho);
    return x;
}
\end{lstlisting}
    \end{column}
  \end{columns}
  \begin{center}
  \begin{tikzpicture}
    \node[var] (i) at (0,0) {ini};
    \foreach \v/\k in {3/1, 4/2, 5/3} {
      \node[cel, minimum width=0.6cm] (v\k) at (0.7+1.3*\k,0) {\v};
      \node[cel, fill=white, minimum width=0.3cm, right=0pt of v\k] (p\k) {};
    }
    \node[var] (f) at (6.3,0) {fim};
    \draw[seta] (i) -- (v1);
    \draw[seta] (p1.center) -- (v2); \draw[seta] (p2.center) -- (v3);
    \draw[UFSGray] (p3.south west) -- (p3.north east);
    \draw[seta] (f) -- (p3);
  \end{tikzpicture}
  \end{center}
  \fonte{Celes et al.\ (2004), cap.~12~\cite{celes2004introducao};
    Feofiloff (2009), cap.~5~\cite{feofiloff2009algoritmos}}
\end{frame}

\begin{frame}{Por que sai do início e entra no fim}
  \begin{columns}[T]
    \begin{column}{0.48\textwidth}
      \small
      A fila precisa de uma ponta para entrar e outra para sair. Na lista
      simples com \texttt{ini} e \texttt{fim}, três das quatro operações nas
      pontas são baratas:

      \medskip
      \begin{tabular}{@{}lcc@{}}
        & início & fim\\
        \midrule
        inserir & barato & barato com \texttt{fim}\\
        remover & barato & $n$ passos\\
      \end{tabular}
    \end{column}
    \begin{column}{0.48\textwidth}
      \small Remover o último exige achar o penúltimo --- percorrer a lista
      toda. Então a remoção fica no início (a frente) e a inserção, no fim.
      Com isso, a lista simples basta; não é preciso \texttt{ant}.

      \begin{block}{Sem \emph{overflow}}
        \small Como na pilha em lista, a fila só enche quando acaba a
        memória; o \emph{underflow} continua possível.
      \end{block}
    \end{column}
  \end{columns}
  \fonte{Celes et al.\ (2004), cap.~12~\cite{celes2004introducao};
    Cormen et al.\ (2009), §10.2, exercício 10.2-3~\cite{cormen2009introduction}}
\end{frame}

\begin{frame}[fragile]{E se esquecer o \texttt{fim}?}
  \begin{columns}[c]
    \begin{column}{0.54\textwidth}
\begin{lstlisting}[style=codC]
int desenfileira(Fila *f) /* esvazia.c */
{
    No *velho = f->ini;
    int x = velho->valor;
    f->ini = velho->prox;
    /* sumiu: if (f->ini == NULL) f->fim = NULL; */
    free(velho);
    return x;
}
/* main: */
enfileira(&f, 1);   enfileira(&f, 2);
printf("sai %d\n", desenfileira(&f));
printf("sai %d\n", desenfileira(&f));
enfileira(&f, 3);
printf("vazia? %s\n", vazia(&f) ? "sim" : "nao");
\end{lstlisting}
    \end{column}
    \begin{column}{0.42\textwidth}
      \begin{block}{Pense antes de passar o slide}
        \small \texttt{enfileira} é a de \texttt{fila\_lista.c}.
        \begin{enumerate}
          \item Depois dos dois \texttt{desenfileira}, quanto valem
            \texttt{ini} e \texttt{fim}?
          \item Ao enfileirar o 3, qual ramo do \texttt{if} é seguido?
          \item O que o programa imprime na última linha?
        \end{enumerate}
      \end{block}
    \end{column}
  \end{columns}
\end{frame}

\begin{frame}[fragile]{Execute e compare}
\begin{lstlisting}[style=shellStyle]
$ ./esvazia
sai 1
sai 2
vazia? sim
$ gcc -g -fsanitize=address esvazia.c -o esvazia-asan && ./esvazia-asan
sai 1
sai 2
==96228==ERROR: AddressSanitizer: heap-use-after-free on address 0x602000000118
WRITE of size 8 at 0x602000000118 thread T0
    #0 0x000104bc5934 in enfileira esvazia.c:23
0x602000000118 is located 8 bytes inside of 16-byte region [0x602000000110,0x602000000120)
freed by thread T0 here:
    #1 0x000104bc5a8d in desenfileira esvazia.c:33
\end{lstlisting}
  \begin{center}
    \scriptsize\textcolor{UFSGray}{A linha 23 é o \texttt{f->fim->prox = novo;}
    de \texttt{enfileira}; a 33, o \texttt{free(velho);}. Relatório resumido.}
  \end{center}
\end{frame}

\begin{frame}{A resposta: \texttt{fim} aponta para um nó que já foi liberado}
  \begin{columns}[c]
    \begin{column}{0.44\textwidth}
      \begin{center}
      \begin{tikzpicture}
        \node[var] (i) at (0,0.9) {ini};
        \node[leg, anchor=west] at (0.5,0.9) {\texttt{NULL}};
        \node[var] (f) at (0,0) {fim};
        \node[celmorta, minimum width=0.6cm] (v) at (1.6,0) {2};
        \node[cel, fill=white, minimum width=0.3cm, right=0pt of v] (p) {};
        \draw[setar] (f) -- (v);
        \node[leg, text=UFSRed] at (1.75,-0.45) {liberado};
        \node[celr, minimum width=0.6cm] (n) at (3.2,0) {3};
        \draw[setar] (p.center) to[out=0, in=180] (n);
      \end{tikzpicture}
      \end{center}
    \end{column}
    \begin{column}{0.52\textwidth}
      \small
      Ao sair o 2, \texttt{ini} virou \texttt{NULL}, mas \texttt{fim}
      continuou apontando para o nó do 2 --- já devolvido pelo
      \texttt{free}. Ao enfileirar o 3, \texttt{fim != NULL}: o programa
      escreve \texttt{fim->prox} em memória liberada e deixa \texttt{ini}
      em \texttt{NULL}.

      \medskip
      O 3 foi alocado e ``enfileirado'', mas ninguém chega até ele: a fila
      se diz vazia e o nó vaza. A remoção do \textbf{último} é o caso
      especial da fila em lista: \texttt{ini} e \texttt{fim} mudam juntos.
    \end{column}
  \end{columns}
  \fonte{Celes et al.\ (2004), cap.~12~\cite{celes2004introducao};
    Serebryany et al.\ (2012)~\cite{serebryany2012asan}}
\end{frame}

\begin{frame}[fragile]{Uma interface, duas implementações}
  \begin{columns}[T]
    \begin{column}{0.5\textwidth}
\begin{lstlisting}[style=codC]
/* fila.h */
#ifndef FILA_H
#define FILA_H

typedef struct fila Fila; /* campos no .c */

Fila *fila_cria(void);
void fila_libera(Fila *f);

void enfileira(Fila *f, int x);
int desenfileira(Fila *f);
int frente(const Fila *f);
int vazia(const Fila *f);
int tamanho(const Fila *f);

#endif
\end{lstlisting}
    \end{column}
    \begin{column}{0.46\textwidth}
      \small \texttt{fila\_vetor.c} é a circular que cresce;
      \texttt{fila\_lista.c}, a lista com \texttt{ini} e \texttt{fim}.
      \texttt{batata\_tad.c} tem a mesma \texttt{roda}, agora com
      \texttt{Fila *}, e roda também $n = 100\,000$ com $k = 1$.
      \begin{block}{Pense antes de passar o slide}
        \small
        \begin{enumerate}
          \item O que muda em \texttt{batata\_tad.c} para trocar o vetor
            pela lista?
          \item As duas versões dão o mesmo sobrevivente para
            $n = 100\,000$?
        \end{enumerate}
      \end{block}
    \end{column}
  \end{columns}
  \fonte{Kernighan e Ritchie (1988), §4.5~\cite{kernighan1988c};
    Celes et al.\ (2004), cap.~9~\cite{celes2004introducao}}
\end{frame}

\begin{frame}[fragile]{Execute e compare}
\begin{lstlisting}[style=shellStyle]
$ gcc -Wall -Wextra batata_tad.c fila_vetor.c -o batata_vetor
$ gcc -Wall -Wextra batata_tad.c fila_lista.c -o batata_lista
$ ./batata_vetor                          $ ./batata_lista
n = 6, k = 0: sai 1 2 3 4 5 | fica 6      n = 6, k = 0: sai 1 2 3 4 5 | fica 6
n = 5, k = 2: sai 3 1 5 2 | fica 4        n = 5, k = 2: sai 3 1 5 2 | fica 4
n = 7, k = 1: sai 2 4 6 1 5 3 | fica 7    n = 7, k = 1: sai 2 4 6 1 5 3 | fica 7
n = 41, k = 2: fica 31                    n = 41, k = 2: fica 31
n = 100000, k = 1: fica 68929             n = 100000, k = 1: fica 68929
\end{lstlisting}
  \begin{columns}[T]
    \begin{column}{0.48\textwidth}
      \small \texttt{batata\_tad.c} não mudou \textbf{nenhuma} linha: só o
      arquivo passado ao \texttt{gcc}. Nenhuma das duas versões tem
      \texttt{MAX}.
    \end{column}
    \begin{column}{0.48\textwidth}
      \small Para $k = 1$ há fórmula: escreva $n = 2^m + l$, com
      $0 \le l < 2^m$; sobra $2l + 1$. Para $100\,000 = 2^{16} + 34\,464$,
      sobra $68\,929$.
    \end{column}
  \end{columns}
  \fonte{Graham, Knuth e Patashnik (1994), §1.3, equação (1.9)~\cite{graham1994concrete}}
\end{frame}

\begin{frame}{Recapitulando: fila numa lista}
  \begin{block}{O que aprendemos}
    \begin{itemize}
      \item a frente é o primeiro nó, o fim é o último: enfileirar é
        inserir depois de \texttt{fim}, desenfileirar é remover o primeiro;
      \item ao contrário da pilha, a fila precisa das \textbf{duas} pontas
        --- e de \texttt{fim} para não percorrer a lista;
      \item quando sai o último, \texttt{fim} também vira \texttt{NULL};
        esquecer isso deixa um ponteiro para memória liberada;
      \item com \texttt{fila.h}, quem usa a fila não sabe se ela é vetor ou
        lista.
    \end{itemize}
  \end{block}
\end{frame}

% ==================================================================
\section{Quanto custa}
% ==================================================================

\begin{frame}{A roda com $n$ crianças e $k = 1$: três filas}
  \begin{columns}[c]
    \begin{column}{0.54\textwidth}
      \begin{center}
      \small
      \begin{tabular}{@{}rrrr@{}}
        \toprule
        $n$ & desloca & circular & lista\\
        \midrule
        25\,000  & 0,053\,s & 0,0004\,s & 0,001\,s\\
        50\,000  & 0,239\,s & 0,0006\,s & 0,002\,s\\
        100\,000 & 1,100\,s & 0,0013\,s & 0,005\,s\\
        200\,000 & 4,217\,s & 0,0026\,s & 0,009\,s\\
        \bottomrule
      \end{tabular}

      \medskip
      \scriptsize \texttt{custo.c}, \texttt{-O2}; as três dão o mesmo
      sobrevivente
      \end{center}
    \end{column}
    \begin{column}{0.42\textwidth}
      \small
      Deslocando, dobrar $n$ \textbf{quadruplica} o tempo: cada saída copia
      a fila inteira, e são $n$ saídas.

      \medskip
      Na circular e na lista, dobrar $n$ dobra o tempo: custo fixo por
      operação. A lista é umas 4 vezes mais lenta que a circular --- um
      \texttt{malloc} e um \texttt{free} por passada da batata.
    \end{column}
  \end{columns}
  \fonte{medido com \texttt{custo.c};
    Bryant e O'Hallaron (2016), §6.2 e §9.9~\cite{bryant2016computer}}
\end{frame}

\begin{frame}{Vetor ou lista?}
  \begin{center}
  \small
  \begin{tabular}{@{}lccc@{}}
    \toprule
    & vetor que desloca & circular que dobra & lista com \texttt{fim}\\
    \midrule
    \texttt{enfileira}           & constante & constante (amortizado) & constante + \texttt{malloc}\\
    \texttt{desenfileira}        & \textbf{tamanho da fila} & constante & constante + \texttt{free}\\
    pode encher?                 & sim, em \texttt{MAX} & só sem memória & só sem memória\\
    caso delicado                & --- & crescer na ordem & \texttt{fim} ao esvaziar\\
    \bottomrule
  \end{tabular}
  \end{center}
  \small Como na pilha, o vetor --- aqui, circular --- costuma ganhar. A lista
  compensa quando cada elemento é grande e já mora num nó, ou quando um
  \texttt{enfileira} nunca pode demorar mais que os outros.
  \fonte{Cormen et al.\ (2009), §10.1 e §17.4~\cite{cormen2009introduction};
    Celes et al.\ (2004), cap.~12~\cite{celes2004introducao}}
\end{frame}

\begin{frame}{Recapitulando: quanto custa}
  \begin{block}{O que aprendemos}
    \begin{itemize}
      \item deslocar todo mundo a cada saída custa o tamanho da fila: o
        total cresce com $n^2$;
      \item na fila circular e na lista, toda operação tem custo fixo;
      \item a circular é a mais rápida na prática: memória contígua e
        crescimentos raros;
      \item a lista paga um \texttt{malloc} por \texttt{enfileira} e 16
        bytes por \texttt{int}.
    \end{itemize}
  \end{block}
\end{frame}

% ==================================================================
\section{Mãos à obra}
% ==================================================================

\begin{frame}{Altere o programa}
  \small
  \begin{block}{Tarefa 1}
    Em \texttt{batata.c}, qual é o menor \texttt{MAX} com que
    \texttt{./batata 41 2} termina? Calcule antes, depois confira. E para
    \texttt{./batata 41 0}?
  \end{block}
  \begin{block}{Tarefa 2}
    Faça a batata passar um número \textbf{diferente} de vezes a cada rodada:
    na rodada $r$ (a partir de 1), ela passa $r$ vezes. Preveja a ordem de
    saída para $n = 5$ antes de executar.
  \end{block}
  \begin{block}{Tarefa 3}
    Em \texttt{batata\_circular.c}, troque \texttt{MAX} por 40 e rode
    \texttt{./batata\_circular 41 2}. Em que momento o programa para? Por
    que não para antes?
  \end{block}
\end{frame}

\begin{frame}{Exercícios (1/3)}
  \footnotesize
  \begin{block}{Exercício 1 --- uma fila que avisa}
    Na fila circular, \texttt{enfileira} devolve 0 se está cheia e
    \texttt{desenfileira(f, \&x)}, 0 se está vazia. Na falha, a fila não muda.
  \end{block}
  \begin{block}{Exercício 2 --- inverter uma fila}
    Com uma pilha auxiliar, inverta a ordem de uma fila usando só as
    operações de \texttt{fila.h} e \texttt{pilha.h}.
  \end{block}
  \begin{block}{Exercício 3 --- escalonador \emph{round-robin}}
    O processo $i$ precisa de $t_i$ unidades de tempo. O da frente recebe
    $q$ unidades e, se não terminou, volta para o fim. Em que ordem e em
    que instante cada um termina?
  \end{block}
  \fonte{Celes et al.\ (2004), cap.~12~\cite{celes2004introducao};
    Tanenbaum e Bos (2015), §2.4.3~\cite{tanenbaum2015modern}}
\end{frame}

\begin{frame}{Exercícios (2/3)}
  \footnotesize
  \begin{block}{Exercício 4 --- fila com duas pilhas}
    Implemente \texttt{fila.h} usando só duas pilhas: uma recebe os que
    entram, a outra entrega os que saem. Quando passar de uma para a outra?
    Por que o custo amortizado é constante?
  \end{block}
  \begin{block}{Exercício 5 --- as últimas $k$ linhas}
    Imprima as últimas $k$ linhas da entrada, como \texttt{tail -n k}, com
    um vetor circular de $k$ cadeias.
  \end{block}
  \begin{block}{Exercício 6 --- fila de duas pontas}
    Um \emph{deque} insere e remove nas duas pontas. Implemente-o num vetor
    circular. Como andar para trás, se em C \texttt{-1 \% cap} é $-1$?
  \end{block}
  \fonte{Cormen et al.\ (2009), §10.1, exercícios 10.1-5 e 10.1-6~\cite{cormen2009introduction};
    Knuth (1997), §2.2.1~\cite{knuth1997art}}
\end{frame}

\begin{frame}{Exercícios (3/3)}
  \footnotesize
  \begin{block}{Exercício 7 --- Josephus sem simulação}
    Compare a simulação com a recorrência $J(1) = 0$,
    $J(n) = (J(n-1) + k + 1) \bmod n$ (posições a partir de 0), escrita
    recursivamente e com um laço. Até que $n$ cada uma vai?
  \end{block}
  \begin{block}{Exercício 8 --- média móvel}
    Leia uma sequência de medições e imprima, a cada uma, a média das últimas
    $k$, em tempo constante por medição: guarde as $k$ num vetor circular e
    mantenha a soma.
  \end{block}
  \begin{block}{Exercício 9 --- o caixa do banco}
    Clientes chegam em instantes dados, cada um com um tempo de atendimento.
    Com um caixa, calcule a espera de cada cliente e a média. E com dois
    caixas e uma fila única?
  \end{block}
  \fonte{Graham, Knuth e Patashnik (1994), §1.3~\cite{graham1994concrete};
    Tenenbaum et al.\ (1995), cap.~4~\cite{tenenbaum1995estruturas}}
\end{frame}

\begin{frame}[fragile]{Desafio}
  \begin{columns}[T]
    \begin{column}{0.6\textwidth}
      \begin{block}{O caminho mais curto no labirinto}
        \small Leia um labirinto (\texttt{\#} parede, \texttt{.} livre,
        \texttt{S} saída, \texttt{E} entrada) e imprima o menor número de
        passos de \texttt{E} até \texttt{S}, andando nas quatro direções:
        \vspace{-0.3em}
        \begin{itemize}
          \item enfileire \texttt{E} com distância 0;
          \item tire uma casa da frente e enfileire cada vizinha livre ainda
            não visitada, com distância + 1;
          \item a primeira vez que \texttt{S} sair da fila, a distância dela
            é a resposta --- ou \texttt{sem caminho}.
        \end{itemize}
      \end{block}
    \end{column}
    \begin{column}{0.36\textwidth}
\begin{lstlisting}[style=shellStyle, language={}]
#########
#E..#...#
#.#.#.#.#
#.#...#S#
#########
\end{lstlisting}
      \centering\scriptsize \texttt{passos: 12}

      \medskip
      \raggedright Por que a fila garante o \emph{menor} caminho --- e uma
      pilha, não?
    \end{column}
  \end{columns}
  \fonte{Moore (1959)~\cite{moore1959shortest};
    Feofiloff (2009), cap.~5~\cite{feofiloff2009algoritmos};
    Cormen et al.\ (2009), §22.2~\cite{cormen2009introduction}}
\end{frame}

\begin{frame}{Fechamento}
  \begin{columns}[T]
    \begin{column}{0.5\textwidth}
      \textbf{O que mudou hoje}\\[0.3em]
      \small Usar as duas pontas, em vez de uma, fez o vetor ``andar''. O
      anel e o \texttt{\% cap} reaproveitam o espaço; a lista precisa de
      \texttt{ini} e de \texttt{fim}.
    \end{column}
    \begin{column}{0.5\textwidth}
      \textbf{O que ficou}\\[0.3em]
      \small Pilha e fila escondem os campos atrás de um \texttt{.h} e
      oferecem só operações. Esse é o desenho de um \emph{tipo abstrato de
      dados} --- o tema da próxima aula.
    \end{column}
  \end{columns}
  \begin{block}{Para levar}
    \begin{itemize}
      \small
      \item fila: entra no fim, sai da frente --- o primeiro a entrar sai
        primeiro;
      \item num vetor: anel com \texttt{\% cap} e um contador \texttt{n};
      \item para crescer: copie na ordem da fila, não na do vetor;
      \item numa lista: sai do primeiro nó, entra depois do último; ao
        esvaziar, \texttt{fim = NULL}.
    \end{itemize}
  \end{block}
  \small\textcolor{UFSGray}{Próxima aula: tipos abstratos de dados ---
  implementação com structs e ponteiros.}
\end{frame}

\section*{Referências}
\begin{frame}[allowframebreaks]{Referências}
  \nocite{*}
  \bibliographystyle{plain}
  \bibliography{referencias}
\end{frame}

\end{document}
